\documentclass[compress,dvipdfmx]{beamer}

\usepackage[T1]{fontenc}
\usepackage[english]{babel}
\usepackage{lmodern}
\usepackage{beramono}          % Bold in verbatim
\usepackage{hyphenat}          % \hyp{} is a breakable dash
\usepackage{xspace}
\usepackage{graphicx}
\usepackage{alltt}
\usepackage{amssymb,amsmath,stmaryrd}
\usepackage{mathpartir}
\usepackage{multicol}
\usepackage{clrscode}
\usepackage{url}

%-*-latex-*-

% Shell prompting
%
\newcommand{\shellprompt}{\textcolor{red}{\texttt{}\$}}
\newcommand{\shellin}[1]{\textsf{\shellprompt{}\ {#1}}}

% OCaml toplevel
%
\newcommand{\camlprompt}{\texttt{\textcolor{blue}{\#}}}
%\newcommand{\semi}{\texttt{\textbf{;;}}}
\newcommand{\topin}[1]{\texttt{\textbf{\#}\ {#1}\textbf{;;}}}
\newcommand{\topout}[1]{\emph{\texttt{\textcolor{green}{#1}}}}
\newcommand{\error}[1]{\emph{\texttt{\textcolor{red}{#1}}}}

% Tokens
%
\newcommand{\tok}[1]{\small \cst{#1}\xspace}

\newcommand{\Teof}{\tok{EOF}}

\newcommand{\Tident}{\tok{IDENT}}
\newcommand{\Tint}{\tok{INT}}
\newcommand{\Tplus}{\tok{PLUS}}
\newcommand{\Tminus}{\tok{MINUS}}
\newcommand{\Ttimes}{\tok{TIMES}}
\newcommand{\Tslash}{\tok{SLASH}}
\newcommand{\Tequal}{\tok{EQUAL}}
\newcommand{\Tarrow}{\tok{ARROW}}
\newcommand{\Tlpar}{\tok{LPAR}}
\newcommand{\Trpar}{\tok{RPAR}}
\newcommand{\Tlet}{\tok{LET}}
\newcommand{\Trec}{\tok{REC}}
\newcommand{\Tin}{\tok{IN}}
\newcommand{\Tfun}{\tok{FUN}}
\newcommand{\Tmatch}{\tok{MATCH}}
\newcommand{\Tif}{\tok{IF}}
\newcommand{\Tthen}{\tok{THEN}}
\newcommand{\Telse}{\tok{ELSE}}
\newcommand{\Ttrue}{\tok{TRUE}}
\newcommand{\Tfalse}{\tok{FALSE}}
\newcommand{\Tand}{\tok{AND}}
\newcommand{\Tor}{\tok{OR}}
\newcommand{\Tnot}{\tok{NOT}}

% ocamllex
%
\newcommand{\Xrule}{\kwd{rule}\xspace}
\newcommand{\Xparse}{\kwd{parse}\xspace}

% Abstract Syntax
%
\newcommand{\clos}[3]{\cst{Clos} \, \lpar{#1},#2,#3\rpar}

% Meta-variable
%
\newcommand{\meta}[1]{\overline{#1}}

% Judgements (evaluation and typing)
%
\newcommand{\ceval}[2]{{#1} \twoheadrightarrow {#2}}
\newcommand{\eval}[3]{{#1} \vdash {#2} \twoheadrightarrow {#3}}
\newcommand{\meval}[3]{\eval{#1}{\src{#2}}{#3}}
\newcommand{\evalb}[2]{\ceval{#1}{#2}}
\newcommand{\ieval}[5]{{#1}/{#2} \vdash {#3} \twoheadrightarrow {#4}/{#5}}
\newcommand{\trel}[3]{{#1} \vdash {#2} : {#3}}
\newcommand{\dom}[1]{\mathrm{dom} (#1)}
\newcommand{\codom}[1]{\mathrm{codom} (#1)}
\newcommand{\subst}[3]{{#1}[#2 \leftarrow #3]}

% Syntax analyser
%
\newcommand{\src}[1]{\langle\!\langle{#1}\rangle\!\rangle}
%\newcommand{\src}[1]{\llbracket{#1}\rrbracket}

% Decorated AST node 
%
\newcommand{\mknode}[3]%
  {#1~[tnpos=l]{\textbf{#2}}~[tnpos=r]{\textcolor{red}{#3}}}

% Comments and examples
%
\newcommand{\remarque}{\textbf{Remarque}\xspace}
\newcommand{\remarques}{\textbf{Remarques}\xspace}
\newcommand{\exemple}{\textbf{Exemple}\xspace}
\newcommand{\exemples}{\textbf{Exemples}\xspace}

% Regular formal langages
%
%\newcommand{\R}[1]{{\cal R} ({#1}^{*})}
\newcommand{\Reg}[1]{{\cal R} ({#1}^{*})}

% AST annotation (typed expression)
%
\newcommand{\te}[2]{{#1}^{#2}}

% Type Unification
%
\newcommand{\mgu}{\textsf{mgu}\xspace}

%%-*-latex-*-


% OCaml pieces of source code
%
\newcommand{\phrase}[1]{\texttt{#1\texttt{\scriptsize ;;}}}
\newcommand{\excerpt}[1]{\texttt{#1}}

% OCaml language
%
\newcommand{\ipat}[1]{\overline{#1}}
\newcommand{\topclos}{\langle{fun}\rangle}

\newcommand{\wild}{{\LARGE \_}}
\newcommand{\kwd}[1]{\textsf{\textbf{#1}}}
\newcommand{\ident}[1]{\textsf{#1}}
\newcommand{\cst}[1]{\textsf{#1}}
\newcommand{\type}[1]{\textsl{#1}}
\newcommand{\str}[1]{\textsf{"#1"}}
\newcommand{\num}[1]{\textsf{#1}}
\newcommand{\comment}[1]{\textsf{(* #1 *)}}
%\newcommand{\cons}{\textsf{::}}
\newcommand{\semi}{\textsf{;;}}
\newcommand{\lpar}{\textsf{(}}
\newcommand{\rpar}{\textsf{)}}
\newcommand{\lbra}{\textsf{[}}
\newcommand{\rbra}{\textsf{]}}
\newcommand{\equal}{\textsf{=}\xspace}
\newcommand{\nequal}{\textsf{\footnotesize <>}\xspace}
\newcommand{\assign}{\textsf{:=}\xspace}
\newcommand{\unit}{\textsf{()}\xspace}
\newcommand{\vbar}{\texttt{|}\xspace}

\newcommand{\Xopen}{\kwd{open}\xspace}
\newcommand{\Xlet}{\kwd{let}\xspace}
\newcommand{\Xrec}{\kwd{rec}\xspace}
\newcommand{\Xin}{\kwd{in}\xspace}
\newcommand{\Xfun}{\kwd{fun}\xspace}
\newcommand{\Xmatch}{\kwd{match}\xspace}
\newcommand{\Xtry}{\kwd{try}\xspace}
\newcommand{\Xwith}{\kwd{with}\xspace}
\newcommand{\Xif}{\kwd{if}\xspace}
\newcommand{\Xthen}{\kwd{then}\xspace}
\newcommand{\Xelse}{\kwd{else}\xspace}
\newcommand{\Xtype}{\kwd{type}\xspace}
\newcommand{\Xand}{\kwd{and}\xspace}
\newcommand{\Xof}{\kwd{of}\xspace}
\newcommand{\Xfor}{\kwd{for}\xspace}
\newcommand{\Xwhile}{\kwd{while}\xspace}
\newcommand{\Xdo}{\kwd{do}\xspace}
\newcommand{\Xdone}{\kwd{done}\xspace}
\newcommand{\Xbegin}{\kwd{begin}\xspace}
\newcommand{\Xend}{\kwd{end}\xspace}
\newcommand{\Xexception}{\kwd{exception}\xspace}
\newcommand{\Xas}{\kwd{as}\xspace}
\newcommand{\Xtrue}{\kwd{true}\xspace}
\newcommand{\Xfalse}{\kwd{false}\xspace}

% Beamer configuration
%
\usetheme{default}
\setbeamersize{text margin left=6mm}
\setbeamersize{text margin right=6mm}
\setbeamertemplate{headline}{}
\setbeamertemplate{navigation symbols}{}
\setbeamertemplate{footline}[page number]{}
\setbeamertemplate{itemize items}[circle]
\setbeamertemplate{theorems}[numbered]


\newcommand\Clang{\textsf{C}\xspace}
\newcommand{\Ada}{\textsf{Ada}\xspace}
\newcommand{\Java}{\textsf{Java}\xspace}
\newcommand{\OCaml}{\textsf{OCaml}\xspace}
\newcommand{\Maude}{\textsf{Maude}\xspace}
\newcommand{\Csharp}{\textsf{C}\raise 0.7mm \hbox{\(\scriptscriptstyle \textsf{\#}\)}\xspace}
\newcommand{\NaN}{\textsf{NaN}\xspace}
\newcommand{\true}{\textsf{true}\xspace}
\newcommand{\false}{\textsf{false}\xspace}
\newcommand\fun[1]{\textsf{#1}}
\newcommand\el{[\,]}
\newcommand\cons[2]{{#1}\!::\!{#2}}

\title{A Theory of Programming}
\author{Christian Rinderknecht}
\institute{Eszterházy Károly Katolikus Egyetem}
\date{September 2026}

\begin{document}

\frame{\maketitle}

% ------------------------------------------------------------------------
%
\begin{frame}{Introduction}

  This lecture proposes a stepwise introduction to a foundation of
  programming and the analyses to write reliable and efficient
  software.

  \bigskip

  There are many different frameworks. We choose here a formal logic
  where relations and functions are the main building blocks. This
  entails that the meaning of \textbf{functional programs} can often
  be made very precise. (We will cover imperative programs later.)

  \bigskip

  It is a good framework to get a sense of \textbf{static analyses},
  and, as such, it stands a prominent place in the field of formal
  methods and \textbf{compiler construction}.

\end{frame}

% ------------------------------------------------------------------------
%
\begin{frame}{Introduction}

  In a later lecture, we will learn a specific programming language
  called \textbf{\OCaml}, enabling you to put in practice what you
  have learnt.

  \bigskip

  Most textbooks and courses on \OCaml propose a traditional approach:
  basic concepts, data structures, algorithms.

  \bigskip

  By contrast, this lecture is not specific to \OCaml and is
  theoretical.

\end{frame}

% ------------------------------------------------------------------------
%
\begin{frame}{String rewriting}
  Let us consider a string of white and black beads, like \(\circ
  \bullet \bullet \bullet \circ \circ \bullet\), and a one-player game
  whose unique move consists in replacing two adjacent beads with only
  one according to the rules
  \begin{equation*}
    \bullet \; \circ   \xrightarrow{\alpha} \bullet\qquad\qquad
    \circ   \; \bullet \xrightarrow{\beta} \bullet\qquad\qquad
    \bullet \; \bullet \xrightarrow{\gamma} \circ
  \end{equation*}

  The rules~\(\alpha\), \(\beta\)~and~\(\gamma\) make up a simple
  \textbf{string\hyp{}rewriting system}.

  \bigskip

  Rules \(\alpha\)~and~\(\beta\) can be conceived as:\\ \emph{`A black
  bead absorbs the white bead next to it.'}

  \bigskip

  The goal of this game is to end up with as few beads as possible, so
  our example may lead to the \textbf{rewrites}
  \begin{equation*}
    \circ \bullet \bullet \, \fbox{\(\bullet \, \circ\)} \circ \bullet
    \xrightarrow{\alpha} \circ \bullet \bullet \, \fbox{\(\bullet \,
      \circ\)} \, \bullet \xrightarrow{\alpha} \fbox{\(\circ \,
      \bullet\)} \bullet \bullet \bullet \xrightarrow{\beta} \bullet
    \bullet \fbox{\(\bullet \, \bullet\)} \xrightarrow{\gamma} \bullet \,
    \fbox{\(\bullet \, \circ\)} \xrightarrow{\alpha} \fbox{\(\bullet
      \, \bullet\)} \xrightarrow{\gamma} \circ
  \end{equation*}

  \medskip

  \noindent where the part of the string to be rewritten next is
  framed.

\end{frame}

% ------------------------------------------------------------------------
%
\begin{frame}{Normal forms and termination}

  Other compositions of the rules lead to the same result~\(\circ\) as
  well. Some others bring all\hyp{}white or all\hyp{}black strings,
  like \(\circ \, \circ\) and~\(\bullet\).

  \bigskip

  Strings that can not be further rewritten, or \textbf{reduced}, are
  called \textbf{normal forms}. We may wonder whether all strings have
  a normal form, and, if so, if it is unique and, furthermore, if it
  is either all\hyp{}white or all\hyp{}black.

  \bigskip

  First, let us note that the system is \textbf{terminating}, that is,
  there is no infinite chain of rewrites, because the number of beads
  strictly decreases in all the rules.

  \bigskip

  This is not a necessary condition in general, for instance, \(\circ
  \; \bullet \xrightarrow{\beta} \bullet \circ \circ\) would preserve
  termination because the composition \(\beta\alpha\alpha\) would be
  equivalent to the original rule~\(\beta\).

  \bigskip

  In particular, this means that any string has a normal form.

\end{frame}

% ------------------------------------------------------------------------
%
\begin{frame}{Invariants}

  Second, notice how the parity of the number of black beads is
  \textbf{invariant} through each rule and how there is no rewrite
  rule for two adjacent white beads.

  \bigskip

  Therefore, if there are \(2p\) initial black beads, then composing
  rules \(\alpha\)~and~\(\beta\) lead to an all\hyp{}black string,
  which can be reduced by applying rule~\(\gamma\) to contiguous pairs
  of beads into an all\hyp{}white string made of \(p\)~beads.

  \bigskip

  Otherwise, the same all\hyp{}black string can be reduced by applying
  alternatively \(\gamma\)~and~\(\beta\) on the left end or
  \(\gamma\)~and~\(\alpha\) on the right end, yielding~\(\circ\).

  \bigskip

  Similarly, if there is an initial odd number of black beads, we
  always end up with one black bead.

\end{frame}

% ------------------------------------------------------------------------
%
\begin{frame}{Confluence}

  Consequently, normal forms are not unique. A simple example is
  \begin{equation*}
    \circ \circ \xleftarrow{\gamma} \bullet \bullet \circ
    \xrightarrow{\alpha} \bullet \bullet \xrightarrow{\gamma}
    \circ
  \end{equation*}

  \medskip

  Systems where normal forms are unique are \textbf{confluent}.

  \bigskip

  If we think of a rewrite step as an elementary computation, it
  becomes clear that confluence is a desirable property for carrying
  out computations, as we would expect, for instance, an arithmetic
  expression to have the same value no matter the order of the small
  steps.

  \bigskip

  A systemic order of rule application is called a \textbf{reduction
    strategy}.

\end{frame}

% ------------------------------------------------------------------------
%
\begin{frame}{Completing a system}

  We can obtain confluence by adding the rule~\(\delta\):
  \begin{equation*}
    \bullet \; \circ   \xrightarrow{\alpha} \bullet\qquad\qquad
    \circ   \; \bullet \xrightarrow{\beta} \bullet\qquad\qquad
    \bullet \; \bullet \xrightarrow{\gamma}
    \circ\qquad\qquad
    \circ\; \circ \xrightarrow{\delta} \circ
  \end{equation*}
  The result of the game is now always one bead whose colour depends
  on the original parity of the black beads as before.

  \bigskip

  To see why, let us consider first that two non\hyp{}overlapping
  parts of a string can be rewritten in parallel, so they are not a
  concern.

  \bigskip

  The interesting cases occur when two applications of rules (maybe of
  the same rule) lead to different strings because they do
  overlap. For instance,
  \begin{equation*}
    \circ \, \circ \xleftarrow{\gamma} \bullet \bullet
    \circ \xrightarrow{\alpha} \bullet \, \bullet
  \end{equation*}

  \medskip

  The important point is that \(\circ \, \circ\) and \(\bullet \,
  \bullet\) can be rewritten into~\(\circ\) at the next step by
  \(\delta\)~and~\(\gamma\), respectively.

\end{frame}

% ------------------------------------------------------------------------
%
\begin{frame}{Critical pairs}

  In general, what matters is that all pairs of strings resulting from
  the application of overlapping rules, called \textbf{critical
    pairs}, can be rewritten to the same string: they are
  \textbf{joinable}.

  \bigskip

  In our example, all interactions occur on substrings made of three
  beads, so we must examine eight cases:
  \begin{center}
    \includegraphics[]{bullets}
  \end{center}

  In all the cases, the divergences are joinable in one step at most.

\end{frame}

% ------------------------------------------------------------------------
%
\begin{frame}{Local confluence}

  If all critical pairs are joinable after a divergence in one ore
  more rewrites, the system is \textbf{locally confluent}.

  \bigskip

  Moreover, if the system is also terminating, then \textbf{every
    string has exactly one normal form}.

  \bigskip

  Sometimes, as seen previously, a non\hyp{}confluent system can be
  completed to become confluent, via a \textbf{completion procedure}.

  \bigskip

  In the case of terminating systems, one of such semi\hyp{}decision
  procedures is the \textbf{Knuth\hyp{}Bendix algorithm}.

\end{frame}

% ------------------------------------------------------------------------
%
\begin{frame}{Ground and linear rewrite systems}

  The system we defined is \textbf{ground}, that is, it involves no
  variables.

  \bigskip

  These allow a finite system to denote an infinite number of ground
  rules if variables range over infinite sets, or simply, they reduce
  the size of rewrite systems.

  \bigskip

  For instance, our continued example is equivalent to
  \begin{equation*}
    \bullet \; \circ \xrightarrow{\alpha} \bullet
    \qquad\qquad
    \circ \; x \xrightarrow{\beta+\delta} x
    \qquad\qquad
    \bullet \; \bullet \xrightarrow{\gamma} \circ
  \end{equation*}

  \medskip

  If we accept multiple occurrences of a variable on the
  left\hyp{}hand side of a rule (that is, not left\hyp{}linear), we
  can further decrease the size of the system:
  \begin{equation*}
    x  \; x \xrightarrow{\gamma+\delta} \circ\qquad\qquad
    x  \; y \xrightarrow{\alpha+\beta} \bullet
  \end{equation*}

\end{frame}

% ------------------------------------------------------------------------
%
\begin{frame}{Ordered rewrite systems}

  Notice that there is now an implicit order over the rules: the rule
  \(\gamma+\delta\) must be examined first for a \textbf{match} with a
  part of the current string, because it is included in the second
  (set \(x=y\) in \(\alpha+\beta\)).

  \bigskip

  Non\hyp{}linear systems are equational theories due the implicit
  equality on substructures, which is complex: we will limit ourselves
  to \textbf{linear systems}.

  \bigskip

  In general, the order in which the rules are written on the page is
  their \textbf{logical order}.

  \bigskip

  Moreover, let us remark that the system does not specify that
  \(x\)~must either be~\((\circ)\) or~\((\bullet)\): this should be
  stated nevertheless somewhere else.

  \bigskip

  Generally speaking, this means that the \textbf{type} of the
  variables has to be defined elsewhere or inferred from the variable
  uses.

\end{frame}

% ------------------------------------------------------------------------
%
\begin{frame}{A simple calculator}

  Perhaps, after playing with beads, the simplest thing is an integer
  calculator, that is, integers and the four basic arithmetic
  operators.

  \bigskip

  An \textbf{expression} may be built using these operations, integers
  and parentheses. For example \((2+3)\times(4+5)\) is an expression.

  \bigskip

  A \textbf{value}  is an integer, therefore  it is a special  case of
  expression.

  \bigskip

  What a calculator does is to match an input expression with a value,
  the result. For example, the above expression evaluates in~\(45\).

  \bigskip

  This process is called \textbf{evaluation}.

\end{frame}

% ------------------------------------------------------------------------
%
\begin{frame}{Evaluation as rewriting expressions / Reductions}

  The evaluation of arithmetic expressions can be defined as a series
  of transformations leading to a value (the result). These
  transformations are called \textbf{rewrites} and a series of
  rewrites is a \textbf{reduction} or, more generally, a
  \textbf{computation}.

  \bigskip

  As we learnt in school long time ago, the way to perform these
  rewrites is to select a part of the current expression consisting of
  an operator and its arguments \emph{as values}, then to replace this
  sub-expression by the operation result.

  \bigskip

  In the following reduction, let us underline the sub-expression to
  be rewritten and print in bold the result of each step:
  \begin{equation*}
    (2+3) \times (\underline{4+5}) \rightarrow (\underline{2+3})
    \times \textbf{9} \rightarrow \underline{\textbf{5} \times 9}
    \rightarrow \textbf{45}
  \end{equation*}

\end{frame}

% ------------------------------------------------------------------------
%
\begin{frame}{Strategy of evaluation and rewrite systems}

  Note that there are different ways to choose the sub-expression. For
  example:
  \begin{equation*}
    (\underline{2+3}) \times (4+5) \rightarrow \textbf{5} \times
    (\underline{4+5}) \rightarrow \underline{5 \times \textbf{9}}
    \rightarrow \textbf{45}
  \end{equation*}
  It is usual to write, as a summary: \((2+3)\times(4+5)
  \xrightarrow{*} 45\) or \((2+3)\times(4+5)=45\).

  \bigskip

  So, there are two concepts involved here:
  \begin{enumerate}

    \item the selection of a sub-expression made of an operator and
      its arguments as values,

    \item the rewriting of such sub-expression into a value.

  \end{enumerate}
  They are defined by a \textbf{rewrite system}.

\end{frame}

% ------------------------------------------------------------------------
%
\begin{frame}{Rewrite rules}

  A rewrite system is a set of \textbf{rewrite rules}.

  \bigskip

  Rewriting rules specify how to perform a rewrite.

  \bigskip

  For example, let us assume that we have an infinite set of rewriting
  rules, that is. an infinite rewrite system, like:
  \begin{align*}
    0 + 0 & \rightarrow 0\\
    0 + 1 & \rightarrow 1\\
    \dots & \rightarrow \dots\\
    124 - 57 & \rightarrow 67\\
    \dots & \rightarrow \dots
  \end{align*}

\end{frame}

% ------------------------------------------------------------------------
%
\begin{frame}{Conditional rewrite rules and meta-variables}

  This rewrite system is not enough for the kind of expressions we
  want to evaluate, where it is possible to combine operations, like
  \((2+3) \times (4+5)\).

  \bigskip

  So let us extend our system with rules that specify the
  sub-expressions we can rewrite.

  \bigskip

  Consider the multiplication:
  \begin{equation*}
    e_1 \times e_2 \rightarrow e'_1 \times e_2 \quad \text{if} \; e_1
    \rightarrow e'_1
  \end{equation*}
  where \(e_1\), \(e'_1\) and \(e_2\) denote any expression and are
  called \textbf{meta\hyp{}variables}.

\end{frame}

% ------------------------------------------------------------------------
%
\begin{frame}{Instance of a rewrite rule}

  An \textbf{instance} of this rule is obtained when the
  meta-variables are replaced by expressions.

  \bigskip

  For example, \( (2+3) \times (4+5) \rightarrow 5 \times (4+5) \) is
  an instance since \(2 + 3 \rightarrow 5\). To understand why,
  replace \(e_1\) by \(2+3\), \(e'_1\) by \(5\) and \(e_2\) by
  \(4+5\).

  \bigskip

  Let us rewrite \(((2 + 3) \times 6) \times (4+5)\). First, we see
  that the previous rule has the required shape, i.e., the left side
  of the arrow is \(e_1 \times e_2\) and if \(e_1 = (2+3) \times 6\)
  and \(e_2 = 4+5\) then \(e_1 \times e_2\) is exactly the expression
  we wish to rewrite. So, we replace \(e_1\) and \(e_2\) by their
  associated expressions in the rule and we get an instance of it:
  \begin{equation*}
    ((2 + 3) \times 6) \times (4+5) \rightarrow e'_1 \times (4+5)
    \quad \text{if} \; (2+3) \times 6 \rightarrow e'_1
  \end{equation*}
  Therefore, we now know that we have to rewrite \((2 + 3) \times 6\).

\end{frame}

% ------------------------------------------------------------------------
%
\begin{frame}{Instance of a rewrite rule}

  We need another instance of the same rule. This time, let \(e_1 =
  2+3\) and \(e_2 = 6\). We get
  \begin{equation*}
    (2+3) \times 6 \rightarrow e'_1 \times 6 \quad \text{if} \; 2 + 3
    \rightarrow e'_1
  \end{equation*}
  It is important to understand that this latter \(e'_1\) is not the
  same as the former: they belong to two different instances of the
  same rule. The latter \(e'_1\) is actually easy to find, since we
  have all the rules for basic cases as \(2 + 3 \rightarrow
  5\). Hence, \(e'_1 = 5\) and the second instance is now complete:
  \begin{equation*}
    (2+3) \times 6 \rightarrow 5 \times 6 \quad \text{if} \; 2 + 3
    \rightarrow 5
  \end{equation*}
  Going backwards, this implies that the former \(e'_1\) is \(5 \times
  6\). So, the first instance is
  \begin{equation*}
    ((2+3) \times 6) \times (4+5) \rightarrow (5 \times 6) \times
    (4+5) \quad \text{if} \; 2 + 3 \rightarrow 5
  \end{equation*}

\end{frame}

% ------------------------------------------------------------------------
%
\begin{frame}{Instance of a rewrite rule}

  Our rule for multiplication, \(e_1 \times e_2 \rightarrow e'_1
  \times e_2\), actually depends on another, \(e_1 \rightarrow
  e'_1\). It is usual to write them vertically:
  \begin{mathpar}
    \inferrule*[right=\(\TirName{Mult}_1\)]
               {e_1 \rightarrow e'_1}
               {e_1 \times e_2 \rightarrow e'_1 \times e_2}
  \end{mathpar}
  This is called an \textbf{inference rule}. Note that it has a name,
  \(\TirName{Mult}_1\). The rewrite rule on the top, acting as a
  condition for the rule on the bottom, is called a
  \textbf{premise}. The rule at the bottom is called a
  \textbf{conclusion}.

  \bigskip

  A \textbf{substitution} is a mapping from meta-variables to
  expressions which applies to an inference rule or rewrite rule. For
  example, let \(\sigma\) be the following substitution: \(\{e_1
  \mapsto (2+3) \times 6; e_2 \mapsto 4+5\}\). This notation is
  equivalent to \(\sigma(e_1) = (2+3) \times 6\) and \(\sigma(e_2) =
  4+5\).

\end{frame}

% ------------------------------------------------------------------------
%
\begin{frame}{Instance of a rewrite rule}

  Now we can make a clear definition of the notion of `instance of a
  rule':

  \medskip

  \begin{quote}
    \emph{An instance of a rule is a rule in which some meta-variables
    \(e_i\) are replaced by expressions \(\sigma(e_i)\).}
  \end{quote}
  For example, here is the instance of \(\TirName{Mult}_1\) using
  previously defined substitution \(\sigma\), and noted
  \(\sigma(\TirName{Mult}_1)\):
  \begin{mathpar}
    \inferrule*[right=\(\sigma(\TirName{Mult}_1)\)]
               {(2 + 3) \times 6 \rightarrow e'_1}
               {((2+3) \times 6) \times (4+5) \rightarrow e'_1 \times (4+5)}
  \end{mathpar}

\end{frame}

% ------------------------------------------------------------------------
%
\begin{frame}{Instance of a rewrite rule}

  Now how do we find (a mapping for) \(e'_1\)?

  \bigskip

  We define another substitution \(\sigma'=\{e_1 \mapsto 2 + 3; e_2
  \mapsto 6\}\) and apply it:
  \begin{mathpar}
    \inferrule*[right=\(\sigma'(\TirName{Mult}_1)\)]
        {2 + 3 \rightarrow e'_1}
        {(2 + 3) \times 6 \rightarrow e'_1 \times 6}
  \end{mathpar}
  Note that the meta-variable \(e'_1\) in
  \(\sigma'(\TirName{Mult}_1)\) is not the same as in
  \(\sigma(\TirName{Mult}_1)\): they should have a different mapping.

  \bigskip

  We have a basic rewrite rule \(2 + 3 \rightarrow 5\), thus we extend
  \(\sigma'\) with \(e'_1 \mapsto 5\) and get
  \begin{mathpar}
    \inferrule*[right=\(\sigma'(\TirName{Mult}_1)\)]
        {2 + 3 \rightarrow 5}
        {(2 + 3) \times 6 \rightarrow 5 \times 6}
  \end{mathpar}

\end{frame}

% ------------------------------------------------------------------------
%
\begin{frame}{Instance of a rewrite rule}

  Therefore we can extend \(\sigma\) with \(e'_1 \mapsto 5 \times 6\)
  and we get
  \begin{mathpar}
    \inferrule*[right=\(\sigma(\TirName{Mult}_1)\)]
       {(2 + 3) \times 6 \rightarrow 5 \times 6}
       {((2+3) \times 6) \times (4+5) \rightarrow (5 \times 6) \times (4+5)}
  \end{mathpar}
  We can summarise this rewrite using the two instances of
  \(\TirName{Mult}_1\) as
  \begin{mathpar}
    \inferrule*[right=\(\sigma(\TirName{Mult}_1)\)] {
      \inferrule*[right=\(\;\sigma'(\TirName{Mult}_1)\)]
                 {2 + 3 \rightarrow 5}
                 {(2 + 3) \times 6 \rightarrow 5 \times 6}
    }
    {((2+3) \times 6) \times (4+5) \rightarrow (5 \times 6) \times (4+5)}
  \end{mathpar}

\end{frame}

% ------------------------------------------------------------------------
%
\begin{frame}{Instance of a rewrite rule}

  But what about the rewrite of an expression like \(7 \times (4+5)\)?

  \bigskip

  Let us define another substitution on \(\TirName{Mult}_1\):
  \(\sigma' = \{e_1 \mapsto 7; e_2 \mapsto 4+5\}\).

  \bigskip

  The corresponding instance is
  \begin{mathpar}
    \inferrule*[right=\(\sigma'(\TirName{Mult}_1)\)]
        {7 \rightarrow e'_1}
        {7 \times (4 + 5) \rightarrow e'_1 \times (4 +5)}
  \end{mathpar}
  But there is no rule of shape `\(v \rightarrow \dots\)', where \(v\)
  is a \emph{value} (because a value is, by definition, an expression
  completely reduced).

\end{frame}

% ------------------------------------------------------------------------
%
\begin{frame}{Completing the multiplication rule}

  This means that we need another rule for the evaluation of the
  \emph{right} argument of \((\times)\):
  \begin{mathpar}
    \inferrule*[right=\(\TirName{Mult}_2\)]
        {e_2 \rightarrow e'_2}
        {e_1 \times e_2 \rightarrow e_1 \times e'_2}
  \end{mathpar}
  This inference rule says that we can rewrite a product by rewriting
  its right argument.

  \bigskip

  For example, here is the instance \(\sigma'(\TirName{Mult}_2)\):
  \begin{mathpar}
    \inferrule*[right=\(\sigma'(\TirName{Mult}_2)\)]
        {4 + 5 \rightarrow 9}
        {7 \times (4+5) \rightarrow 7 \times 9}
  \end{mathpar}

\end{frame}

% ------------------------------------------------------------------------
%
\begin{frame}{Strategy}

  But now, something interesting happens.

  \bigskip

  Consider again \((2 + 3) \times (4+5)\): both \(\TirName{Mult}_1\)
  and \(\TirName{Mult}_2\) can be instantiated by substitution
  \(\sigma'' = \{e_1 \mapsto 2+3; e_2 \mapsto 4+5\}\):
  \begin{mathpar}
    \inferrule*[right=\(\sigma''(\TirName{Mult}_1)\)]
        {2 + 3 \rightarrow 5}
        {(2+3) \times (4+5) \rightarrow 5 \times (4+5)}
    \and
    \inferrule*[right=\(\sigma''(\TirName{Mult}_2)\)]
        {4 +5 \rightarrow 9}
        {(2+3) \times (4+5) \rightarrow (2+3) \times 9}
  \end{mathpar}

\end{frame}

% ------------------------------------------------------------------------
%
\begin{frame}{Strategy}

  This means that \emph{our system does not specify the order of
  evaluation of the arguments to \((\times)\).}

  \bigskip

  In other words, our rewrite system does not select only one
  sub\hyp{}expression to rewrite at each step.

  \bigskip

  We need a \textbf{strategy} to complete it.

\end{frame}

% ------------------------------------------------------------------------
%
\begin{frame}{Finiteness and soundness of reductions}

  At this point, it is clear that these reductions satisfy the
  following properties:
  \begin{itemize}

    \item \textbf{Finiteness.} All the reductions are finite, i.e.,
      all the computations terminate.

    \item \textbf{Soundness.} All the reductions of an expression end
      with the same value, if any. (Think confluence.)

  \end{itemize}

  About the second point, note indeed that \emph{some expressions have
  no value}, as
  \begin{equation*}
    (\underline{2+3})/(4-4) \rightarrow \textbf{5}/(\underline{4-4})
    \rightarrow 5/\textbf{0} \nrightarrow
  \end{equation*}
  The reduction is finite but does not end with a value: this
  situation corresponds to an \textbf{error}, usually called
  \textbf{run\hyp{}time error} in programming languages. Expression
  \(5/0\) is \textbf{stuck}.

\end{frame}

% ------------------------------------------------------------------------
%
\begin{frame}{Another model of run-time errors}

  We can modify the calculator so that, in case of error (here, the
  only one is the division by zero), the result of a rewrite is a
  special value, called \NaN (\emph{Not a Number}). Then we add the
  following rules:
  \begin{mathpar}
    \inferrule*[right=\(\quad\TirName{DivZero}\)]
       {}
       {e/0 \rightarrow \NaN}
    \and
    \inferrule*[right=\(\quad\TirName{Div-Err}_1\)]
       {}
       {\NaN/e \rightarrow \NaN}
    \and
    \inferrule*[right=\(\quad\TirName{Div-Err}_2\)]
       {}
       {e/\NaN \rightarrow \NaN}
    \and
    \inferrule*[right=\(\quad\TirName{Mult-Err}_1\)]
      {}
      {\NaN \times e \rightarrow \NaN}
    \and
    \inferrule*[right=\(\quad\TirName{Mult-Err}_2\)]
       {}
       {e \times \NaN \rightarrow \NaN}
    \and
    \inferrule*[right=\(\quad\TirName{Add-Err}_1\)]
       {}
       {\NaN + e \rightarrow \NaN}
    \and
    \inferrule*[right=\(\quad\TirName{Add-Err}_2\)]
       {}
       {e + \NaN \rightarrow \NaN}
    \and
    \inferrule*[right=\(\quad\TirName{Sub-Err}_1\)]
       {}
       {\NaN - e \rightarrow \NaN}
    \and
    \inferrule*[right=\(\quad\TirName{Sub-Err}_2\)]
       {}
       {e - \NaN \rightarrow \NaN}
\end{mathpar}

\end{frame}

% ------------------------------------------------------------------------
%
\begin{frame}{Another model of run-time errors}

  The rationale of this system is that once an error occurs, no other
  computations (leading to integer values) are required and the
  reduction can end as soon as possible.

  \bigskip

  In this case:
  \begin{equation*}
    (2+3)/(4-4) \rightarrow 5/(4-4) \rightarrow 5/0 \rightarrow \NaN
  \end{equation*}
  or, even the shortest possible:
  \begin{equation*}
    (2+3)/(4-4) \rightarrow (2+3)/0 \rightarrow \NaN
  \end{equation*}

\end{frame}

% ------------------------------------------------------------------------
%
\begin{frame}{Rephrasing soundness}

  The interesting phenomenon here is that, now, any reduction ends in
  a value (since errors are values). In the previous rewrite system,
  any reduction ends either in a value or a \textbf{stuck expression},
  i.e., an expression that cannot be rewritten further.

  \bigskip

  With the new system, that explicitly specifies the reduction of the
  \NaN error, the soundness property can be rephrased as

  \medskip

  \begin{quote}
    \emph{All the reductions of an expression end with the same value.}
  \end{quote}
  Obviously, the drawback of the new system is that it requires a lot
  of additional rules, and, as the system is augmented with new rules,
  new rules for the errors have to be added... or are forgotten by
  mistake.

\end{frame}

% ------------------------------------------------------------------------
%
\begin{frame}{Specifying the strategy}

  It is possible to force an order of evaluation on the arguments of
  \((\times)\). Assume we want to reduce the left argument first.

  \bigskip

  The idea is to change \(\TirName{Mult}_2\) so that instead
  of~\(e_1\), which stands for any expression, we specify a
  value~\(v\):
  \begin{mathpar}
    \inferrule*[right=\(\TirName{New-Mult}_2\)]
       {e_2 \rightarrow e'_2}
       {v \times e_2 \rightarrow v \times e'_2}
  \end{mathpar}
  This way, it is not possible anymore to instantiate this rule with
  expression \((2+3) \times (4+5)\) because \(2+3\) is not a value.

\end{frame}

% ------------------------------------------------------------------------
%
\begin{frame}{Specifying the strategy}

  The only possible reduction is:
  \begin{mathpar}
    \inferrule*[right=\(\sigma_1(\TirName{Mult}_1)\)]
       {2 + 3 \rightarrow 5}
       {(2+3) \times (4+5) \rightarrow 5 \times (4+5)}
    \and
    \inferrule*[right=\(\sigma_2(\TirName{New-Mult}_2)\)]
       {4 + 5 \rightarrow 9}
       {5 \times (4+5) \rightarrow 5 \times 9}
  \end{mathpar}
  where \(\sigma_1=\{e_1 \mapsto 2+3; e_2 \mapsto 4+5\}\) and
  \(\sigma_2=\{v \mapsto 5; e_2 \mapsto 4+5\}\).

\end{frame}

% ------------------------------------------------------------------------
%
\begin{frame}{Determinism}

  We have specified the strategy \textbf{in} the rewrite system.

  \bigskip

  Also, with \(\TirName{New-Mult}_2\), given an expression, there is
  only one rule that can be applied (or none): this is sometimes
  called \textbf{determinism}.

  \bigskip

  \textbf{Determinism.} If \(e_1 \rightarrow e'_1\) and \(e_2
  \rightarrow e'_2\) then \(e'_1 = e'_2\).

  \bigskip

  This is a stronger property than soundness, i.e., it implies
  soundness. Indeed, determinism implies that there is only one
  reduction from an expression to its value (if any).

\end{frame}

% ------------------------------------------------------------------------
%
\begin{frame}{Pragmatics of programming languages}

  In many programming languages, the order of evaluation of operator
  arguments is not specified, which usually means that there is a
  sentence in the official document which defines the language like:
  `The order of evaluation of the arguments of operators and functions
  is not specified.'

  \bigskip

  Formally, this means that the underlying rewrite system is
  \textbf{non-deterministic}.

  \bigskip

  One notable exception is \Java, which specifies that the arguments
  are always evaluated from left to right, i.e., following the order
  of writing.

  \bigskip

  Of course, if an order is specified, as in \Java, the compilers of
  these languages must implement it. Otherwise one is chosen
  arbitrarily and may change from one release to another.

\end{frame}

% ------------------------------------------------------------------------
%
\begin{frame}{Rewrite systems as models of languages}

  By studying rewrite systems, we can learn the basic concepts
  involved in programming languages, because rewrite systems are
  extremely simple.

  \bigskip

  There are some languages, like \Maude, whose programs actually are
  rewrite systems.

  \bigskip

  Some compilers transform the input program (called \textbf{source})
  and output a program, called \textbf{byte-code}, which can be
  reduced by a rewrite system.

\end{frame}

% ------------------------------------------------------------------------
%
\begin{frame}{Rewrite systems as models of languages}

  This is the case of \Java. This is why, in general, we need a
  \textbf{\Java virtual machine} (JVM) to run the program: this JVM
  performs the reductions according to a rewrite system.

  \bigskip

  Even if there is no byte-code, we always can interpret what the
  program does by studying a rewrite system which models it.

\end{frame}

% ------------------------------------------------------------------------
%
\begin{frame}{Sequentiality and parallelism}

  Some programming languages, like \Ada, allow the programmer to
  specify that different parts of its program can actually be executed
  separately: this is called, in general, \textbf{parallelism}. This
  parallelism is possible only if the underlying rewrite system is
  non-deterministic.

  \bigskip

  Indeed, if all the strategies lead to the same value (soundness
  property), it is possible to reduce different parts of an expression
  in parallel if the rewrite system allows it.

\end{frame}

% ------------------------------------------------------------------------
%
\begin{frame}{Sequentiality and parallelism}

  For example, the first argument of \((\times)\) can be computed
  \emph{at the same time} as its second argument. In the case of
  \((2+3) \times (4+5)\), this means that we can reduce in parallel
  \(2+3\), with \(\TirName{Mult}_1\), and \(4+5\), with
  \(\TirName{Mult}_2\), then \(5 \times 9\) is finally rewritten in
  \(45\) (after \textbf{synchronisation} of the two reductions).

\end{frame}

%------------------------------------------------------------------------
%
\begin{frame}[fragile]{Order of evaluation in C}

  In case of C, the order of evaluation arguments of arithmetic
  operators and functions is not specified, it actually depends on the
  compiler and its version. Therefore it would be dangerous to rely on
  the order of evaluation in C.

  \bigskip

  Why is the order left unspecified? Simply because of possible
  \textbf{side-effects}. Imagine the following situation:
{\small
\begin{verbatim}
int n = 0;
int f () { n++; return n; };
int g () { n = 2; return n; };
int main () { return f() + g() + n; }
\end{verbatim}
}
What is the result?

\end{frame}

% ------------------------------------------------------------------------
%
\begin{frame}{Boolean expressions}

  Let~\(\wedge\) represent the conjunction (`and'), \(\vee\) the
  disjunction (`or'), \(\neg\) the negation (`not'), and let the
  two values \true and \false.

  \bigskip

  The definition of the conjunction of values is very simple:
  \begin{mathpar}
    \inferrule
       {}
       {\true \wedge \true \rightarrow \true}
    \and
    \inferrule
       {}
       {\true \wedge \false \rightarrow \false}\\
    \inferrule
       {}
       {\false \wedge \true \rightarrow \false}
    \and
    \inferrule
       {}
       {\false \wedge \false \rightarrow \false}
\end{mathpar}

\end{frame}

% ------------------------------------------------------------------------
%
\begin{frame}{Boolean expressions made simpler}

  This system can even be made simpler, by means of meta-variables:
  \begin{mathpar}
    \inferrule*[right=\(\quad\TirName{True}\)]
       {}
       {\true \wedge \true \rightarrow \true}\\
    \inferrule*[right=\(\quad\TirName{False}_1\)]
       {}
       {\false \wedge v \rightarrow \false}\\
    \inferrule*[right=\(\quad\TirName{False}_2\)]
       {}
       {v \wedge \false \rightarrow \false}
  \end{mathpar}
  where \(v\) denotes any value, i.e., either \true or \false.

\end{frame}

% ------------------------------------------------------------------------
%
\begin{frame}{Boolean expressions}

  So, if we define the rewrite of boolean expressions (not just the
  values \true and \false) as
  \begin{mathpar}
    \inferrule*[right=\(\TirName{\(\wedge_1\)}\)]
       {e_1 \rightarrow e'_1}
       {e_1 \wedge\, e_2 \rightarrow e'_1 \wedge\, e_2}
    \and
    \inferrule*[right=\(\TirName{\(\wedge_2\)}\)]
       {e_2 \rightarrow e'_2}
       {e_1 \wedge\, e_2 \rightarrow e_1 \wedge\, e'_2}
  \end{mathpar}
  then, for example,
  \begin{mathpar}
    \inferrule*[right=\(\sigma_1(\wedge_1)\)]
       {\false \wedge\; \true \rightarrow \false
        \quad \TirName{False}_1}
       {(\false \wedge\; \true) \wedge\, (\true \wedge\, \true)
         \rightarrow \false \wedge\, (\true \wedge\, \true)}
  \end{mathpar}
  where \(\sigma_1=\{e_1 \mapsto \false \wedge \true; e_2 \mapsto
  \true \wedge \true; e'_1 \mapsto \false\}\).

\end{frame}

% ------------------------------------------------------------------------
%
\begin{frame}{Boolean expressions}

  Then
  \begin{mathpar}
    \inferrule*[right=\(\sigma_2(\wedge_2)\)]
       {\true \wedge\, \true \rightarrow \true \quad \TirName{True}}
       {\false \wedge\, (\true \wedge\, \true) \rightarrow \false \wedge\,
  \true}
  \end{mathpar}
  where \(\sigma_2 = \{e_1 \mapsto \false; e_2 \mapsto \true \wedge
  \true; e'_1 \mapsto \true\}\).

  \bigskip

  and, finally, we can instantiate \(\TirName{False}_1\) with
  \(\sigma_3 = \{v \mapsto \true\}\):
  \begin{mathpar}
    \inferrule*[right=\(\quad\sigma_3(\TirName{False}_1)\)]
       {}
       {\false \wedge\, \true \rightarrow \false}
  \end{mathpar}
  In this example, we computed both arguments of \(\wedge\) to reach
  the value.

\end{frame}

% ------------------------------------------------------------------------
%
\begin{frame}{Boolean expressions}

  But the inspection of the basic rules shows that if one of the
  arguments of \(\wedge\) reduces to \false, then there is no need to
  reduce the other argument: the rewrite will always be
  \false. Therefore, it is more efficient to extend the rules on
  values to the expressions:
  \begin{mathpar}
    \inferrule*[right=\(\quad\TirName{False}_1\)]
       {}
       {\false \wedge e \rightarrow \false}
    \and
    \inferrule*[right=\(\quad\TirName{False}_2\)]
       {}
       {e \wedge \false \rightarrow \false}\\
    \inferrule*[right=\(\quad\TirName{True}_1\)]
       {}
       {\true \wedge e \rightarrow e}
    \and
    \inferrule*[right=\(\quad\TirName{True}_2\)]
       {}
       {e \wedge \true \rightarrow e}
  \end{mathpar}
  where \(e\) represents any boolean expression. Then, now
  \begin{mathpar}
    \inferrule*[right=\(\quad\sigma_3(\TirName{False}_1)\)]
       {}
       {\false \wedge\, (\true \wedge\, \true) \rightarrow \false}
  \end{mathpar}
  where \(\sigma_3 = \{e \mapsto \true \wedge \true\}\).

\end{frame}

% ------------------------------------------------------------------------
%
\begin{frame}{Boolean expressions}

  Try to reduce
  \begin{equation*}
    (\true \wedge (\false \wedge \true)) \wedge (\true \wedge (\true
    \wedge \true))
  \end{equation*}
  Note that this rewrite system is not deterministic. If we enforce
  that the left argument of \(\wedge\) is completely reduced first,
  then we can check if it is \true or \false: in the latter case, we
  do not need to reduce the right argument.

\end{frame}

% ------------------------------------------------------------------------
%
\begin{frame}[fragile]{Boolean expressions}

  This is easy to specify:
  \begin{mathpar}
    \inferrule*[right=\(\quad\wedge_{\TirName{True}}\)]
       {}
       {\true \wedge e \rightarrow e}
    \and
    \inferrule*[right=\(\quad\wedge_{\TirName{False}}\)]
       {}
       {\false \wedge e \rightarrow \false}\\
    \inferrule*[right=\(\wedge\)]
       {e_1 \rightarrow e'_1}
       {e_1 \wedge\, e_2 \rightarrow e'_1 \wedge\, e_2}
  \end{mathpar}
  Many programming languages, as C, use this strategy. It allows the
  programmers to write handy code like
{\small
\begin{verbatim}
while (index < max && tab[index] > 0) ...
\end{verbatim}
}

\end{frame}

% ------------------------------------------------------------------------
%
\begin{frame}{Boolean expressions}

  Note that this strategy (left argument before the right one) is not
  optimal in general: the first argument may reduce to
  \true. \emph{Finding an optimal strategy is not always possible.}
  For example, we could add another rule that shortens sometimes the
  reductions:
  \begin{mathpar}
    \inferrule*[right=\(\wedge_{\TirName{Refl}}\)]
       {}
       {e \wedge e \rightarrow e}
  \end{mathpar}
  This rule explicitly specifies the \textbf{reflexivity} of the
  conjunction.

  But the system becomes non-deterministic because there is now an
  overlap between rule \(\wedge_{\TirName{Refl}}\) and both
  \(\wedge_{\proc{True}}\) and \(\wedge_{\proc{False}}\) (consider for
  example \(\true \wedge \false\)).

\end{frame}

% ------------------------------------------------------------------------
%
\begin{frame}{Determinism versus non-determinism}

  Let us summarise the concepts of determinism and
  non\hyp{}determinism.
  \begin{itemize}

    \item Determinism means that, for any expression, there is only
      one rule to rewrite it. This implies that, every time a program
      is run with the same input, the output will be the same. This is
      a useful property, called soundness, in particular for
      debugging.

    \item Non-determinism means that, for any expression, there may be
      several rules to rewrite it. Moreover,
      \begin{itemize}

         \item if the rewrite system is sound, the output will always
           be the same for any run on the same input, even if the
           trace may be different each time, which makes the debugging
           more difficult, but enable parallel reductions.

         \item the rewrite system may be unsound.

      \end{itemize}
  \end{itemize}

\end{frame}

% ------------------------------------------------------------------------
%
\begin{frame}[fragile]{Determinism versus non-determinism}

  Some programming languages have features that may be both
  non\hyp{}deterministic and unsound.

  \bigskip

  Perhaps the most infamous case is in C, where variables do not need
  to be initialised at their declaration:
{\small
\begin{verbatim}
int a;
\end{verbatim}
}
  If the programmer forgets to give a value to the variable, any access
  to it is unspecified.

  \bigskip

  Practically, this means that whatever is located in the memory at
  the location of the variable will be used at each read access. Of
  course, this value may change from one run of the program to
  another, even on the same input.

\end{frame}

% ------------------------------------------------------------------------
%
\begin{frame}{Commutativity and its consequences}

  We may try another rewrite system for \(\wedge\):
  \begin{mathpar}
    \inferrule*[right=\(\quad\TirName{True}\)]
       {\true \wedge e \rightarrow e}
       {}
    \and
    \inferrule*[right=\(\quad\TirName{False}\)]
       {}
       {\false \wedge e \rightarrow e}\\
    \inferrule*[right=\(\wedge\)]
       {e_1 \rightarrow e'_1}
       {e_1 \wedge e_2 \rightarrow e'_1 \wedge e_2}
    \and
    \inferrule*[right=\(\TirName{Comm}\)]
       {}
       {e_1 \wedge e_2 \rightarrow e_2 \wedge e_1}
  \end{mathpar}
  While this system is in theory fine, there are two practical
  problems:
  \begin{enumerate}

    \item it is \textbf{not deterministic} due to \(\TirName{Comm}\)
      that overlaps with all the other rules;

    \item it may \textbf{not terminate} due to rule
      \(\TirName{Comm}\):
      \begin{align*}
        \true \wedge \false
        & \;\xrightarrow{\TirName{Comm}}\;
        \false \wedge \true \;\xrightarrow{\TirName{Comm}}\; \true
        \wedge \false\\
        & \;\xrightarrow{\TirName{Comm}}\; \false
        \wedge \true \rightarrow \cdots
      \end{align*}

  \end{enumerate}

\end{frame}

% ------------------------------------------------------------------------
%
\begin{frame}[fragile]{Non-termination}

  Now some reductions may not be finite, which models
  \textbf{non\hyp{}termination}.

  \bigskip

  This particular kind of non-termination, due to commutativity, is
  the model in programming languages of so-called `\textbf{infinite
    loops}' that do not allocate memory, i.e., that do not require
  additional memory at each loop step. Equivalently, it corresponds to
  a \emph{recursive call} like
{\small
\begin{verbatim}
bool and(bool e1, bool e2) {
  ... return and(e2, e1); ... }
\end{verbatim}
}
  Anyway, in the presence of infinite reductions, the soundness
  property must be weakened as follows:

  \bigskip

  \begin{quote}
    \textbf{Weak soundness.} \emph{All finite reductions of an
    expression end with the same value, if any.}
  \end{quote}

\end{frame}

% ------------------------------------------------------------------------
%
\begin{frame}{Non-termination}

  Let us define the negation:
  \begin{mathpar}
    \inferrule
       {}
       {\neg \true \rightarrow \false}
    \and
    \inferrule
       {}
       {\neg \false \rightarrow \true}
  \end{mathpar}
  It is also well known that the negation satisfies \(e = \neg\neg
  e\). So we could imagine another rewrite rule
  \begin{equation*}
    \neg\neg e \rightarrow e
  \end{equation*}
  This is fine, this rule simplifies double negations. But what if we
  had the converse rule?
  \begin{equation*}
    e \rightarrow \neg\neg e \qquad \text{where} \; e \; \text{is not a
  value.}
  \end{equation*}
  This allows infinite reductions: \(e \rightarrow \neg\neg e
  \rightarrow e \rightarrow \neg\neg e \rightarrow \dots\)

% excluded middle: e or ~e
% non-contradiction: ~(e and ~e)
% de Morgan applied to non-contradiction: ~(e and ~e) = ~e or ~~e

\end{frame}

% ------------------------------------------------------------------------
%
\begin{frame}{Summary}

  There are three kinds of reductions:
  \begin{enumerate}

    \item Finite and ending in a value
  \[e_1 \rightarrow e_2 \rightarrow \dots \rightarrow v\]

    \item Finite but ending in a stuck expression
  \[e_1 \rightarrow e_2 \rightarrow \dots \rightarrow e_n
    \nrightarrow\]

    \item Infinite
  \[e_1 \rightarrow e_2 \rightarrow \cdots\]

  \end{enumerate}

\end{frame}

% ------------------------------------------------------------------------
%
\begin{frame}{Conditional expressions}

  At this point, boolean expressions can be
  \begin{itemize}

    \item \Xtrue (a value)

    \item \Xfalse (a value)

    \item \(e_1 \wedge e_2\)

    \item \(e_1 \vee e_2\)

    \item \(\neg e_1\)

  \end{itemize}
  where \(e_1\) and \(e_2\) are boolean expressions. Let us extend our
  boolean expressions with a \textbf{conditional expression}:
  \begin{itemize}

     \item \Xif \(e_1\) \Xthen \(e_2\) \Xelse \(e_3\)

  \end{itemize}

\end{frame}

% ------------------------------------------------------------------------
%
\begin{frame}{Conditional expressions}

  It is easy to extend the rewrite system to cope with conditionals:
  \begin{mathpar}
    \inferrule*[right=\(\quad\TirName{If-True}\)]
       {}
       {\Xif \; \Xtrue \; \Xthen \; e_2 \; \Xelse \; e_3 \rightarrow e_2}
    \and
    \inferrule*[right=\(\quad\TirName{If-False}\)]
       {}
       {\Xif \; \Xfalse \; \Xthen \; e_2 \; \Xelse \; e_3 \rightarrow e_3}
    \and
    \inferrule*[right=\(\TirName{If}\)]
       {e_1 \rightarrow e'_1}
       {\Xif \; e_1 \; \Xthen \; e_2 \; \Xelse \; e_3
        \rightarrow \Xif \;
        e'_1 \; \Xthen \; e_2 \; \Xelse \; e_3}
  \end{mathpar}
  This means that the boolean expression \(e_1\) must be completely
  reduced \emph{before} the others. This is the usual way in
  programming languages, the rationale being to avoid unnecessary
  computations (evaluating \(e_2\) and~\(e_3\) is not necessary).

\end{frame}

% ------------------------------------------------------------------------
%
\begin{frame}{Term rewriting}

  More general are the \textbf{term\hyp{}rewriting systems}, where a
  \textbf{term} is a mathematical object built on tuples, integers and
  variables.

  \bigskip

   Let us consider the following totally ordered system:
   \begin{align*}
     (0,m) &\xrightarrow{1} m\\
     (n,m) &\xrightarrow{\smash 2} (n-1, n \times m)\\
     n     &\xrightarrow{\smash 3} (n,1)
   \end{align*}

   Arithmetic operators \((-)\)~and~\((\cdot)\) are defined out of the
   system and \(m\)~and~\(n\) are variables denoting natural numbers.

   \bigskip

   Would the rules not be ordered as they are actually laid out, the
   second rule would match any pair. Instead, it can be assumed that
   \(n \neq 0\) when matching with it.

   \bigskip

   The last rule is to be tried last, as it matches all terms.

\end{frame}

% ------------------------------------------------------------------------
%
\begin{frame}{Transitive closures}

  We can easily see that all the compositions of rewrites starting
  with a natural number~\(n\) end with the factorial of~\(n\):
  \smallskip
  \begin{equation*}
    n \xrightarrow{3} (n,1)
    \xrightarrow{2} \dots
    \xrightarrow{2} (0,n!)
    \xrightarrow{1}
    n!, \quad \text{for \(n \in \mathbb{N}\)}.
  \end{equation*}

  Let us note \((\xrightarrow{n})\) the composition of
  \((\rightarrow)\) repeated \(n-1\)~times:
  \begin{align*}
    (\xrightarrow{1})   &:= (\rightarrow)\\
    (\xrightarrow{\smash{n+1}}) & :=
    (\rightarrow) \circ (\xrightarrow{\smash{n}}),
    \quad\text{with \(n > 0\)}.
  \end{align*}

  The \textbf{transitive closure} of \((\rightarrow)\) is defined as
  \begin{equation*}
    (\twoheadrightarrow) := \bigcup_{i > 0}{(\xrightarrow{i})}.
  \end{equation*}

  The factorial function coincides with the transitive closure of
  \((\rightarrow)\): \(n \twoheadrightarrow n!\). Let
  \((\xrightarrow{*})\) be the reflexive\hyp{}transitive closure of
  \((\rightarrow)\), that is,
  \begin{equation*}
    (\xrightarrow{*}) := (=) \cup (\twoheadrightarrow).
  \end{equation*}

\end{frame}

% ------------------------------------------------------------------------
%
\begin{frame}{Functions}

  \begin{itemize}

    \item A confluent system defines a \textbf{partial function}.

    \item If the system is also terminating, we get a
      \textbf{function}.

    \item We can name functions: for example, \(\fun{c}(1,
      \fun{d}(n))\)~is a term constructed with \textbf{function names}
      \fun{c}~and~\fun{d}, as well as variable~\(n\).

    \item A tuple tagged with a function name, like \(\fun{f}(x,y)\),
      is called a \textbf{function call}.

    \item The components of the tuples are then called
      \textbf{arguments}, for example \(\fun{d}(n)\) is the second
      argument of the call \(\fun{c}(1, \fun{d}(n))\).

    \item It is possible for a function call to hold no arguments,
      like \(\fun{d}()\).

    \item We restrict the left\hyp{}hand sides of rules to be function
      calls.

  \end{itemize}

\end{frame}

% ------------------------------------------------------------------------
%
\begin{frame}{Trees}

  The topological understanding of a function call or a tuple is the
  finite \textbf{tree}. A tree is a hierarchical layout of
  information, for example:
  \begin{center}
    \includegraphics{tree_for_term}
  \end{center}

   The disks are called \textbf{nodes} and the segments which connect
   two nodes are called \textbf{edges}. The topmost node (with a
   diameter) is called the \textbf{root} and the bottommost ones
   (\(\bullet\)) are called the \textbf{leaves}. All nodes except the
   leaves are seen to downwardly connect to some other nodes, called
   \textbf{children}. Upwardly, each node but the root is connected to
   another node, called its \textbf{parent}.

\end{frame}

% ------------------------------------------------------------------------
%
\begin{frame}{Trees}

  \textbf{Trees are pervasive} in informatics, especially in
  functional programming, but not only. Structured documents, like
  books, articles, and web pages, are composed of chapters, sections,
  paragraphs, figures, appendices, indices, etc.

  \bigskip

  The occurrences of these components are mutually constrained; for
  instance, it is understood that a section is part of a chapter and
  that appendices are located at the end of a document.

  \bigskip

  This hierarchical layout is meant to facilitate reading, and it
  supports the search for specific items of information.

  \bigskip

  When considering computer systems, these data must be uniformly
  encoded by means of a formal language.

\end{frame}

% ------------------------------------------------------------------------
%
\begin{frame}{Terms as trees}

  Trees can be used to depict terms as follows. A function call is a
  tree whose root is the function name and the children are the trees
  denoting the arguments.

  \bigskip

  A tuple can be considered as having an invisible function name
  represented by a node with a period (\texttt{.}) in the tree, in
  which case the components of the tuple are its children. For
  example,
  \begin{center}
    \includegraphics{tree}
  \end{center}
  is the tree corresponding to the tern
  \(\fun{f}((\fun{g}(0),(x,1)),(),\fun{g}(y))\).

  \bigskip

  Note that tuples, except the empty tuple, are not represented in the
  tree, since they encode the structure itself, which is already laid
  out. The number of arguments of a function is called \textbf{arity}.

\end{frame}

% ------------------------------------------------------------------------
%
\begin{frame}{Functional languages}

  \begin{itemize}

    \item For programming, we only consider confluent systems because
      they define partial functions.

    \item To have only one reduction chain per term to the normal
      form, we may also enforce that rules are ordered.

    \item We also enforce that normal forms must be \textbf{values},
      that is, terms without function calls or only with calls to
      functions not occurring on the left\hyp{}hand sides.

    \item In other words, irreducible calls of defined functions are
      not allowed in values.

    \item Systems with these constraints make up a \textbf{functional
      language}.

    \item A program here is a term, and evaluating it means to rewrite
      the term
      \begin{enumerate}
        \item to a value,
        \item or to end in error with a normal form that is not a value,
        \item or to never terminate.
      \end{enumerate}

  \end{itemize}

\end{frame}

% ------------------------------------------------------------------------
%
\begin{frame}{Call-by-value}

  If a function call contains at least one other function call, there
  is a choice as to which one to reduce first.

  \bigskip

  One reduction strategy, named \textbf{call-by-value}, consists in
  always reducing the arguments before the call itself. Unfortunately,
  call-by-value enables otherwise terminating rewrites to not
  terminate. For instance, let us consider
  \begin{equation*}
    \fun{f}(x) \xrightarrow{\alpha} 0\qquad
    \fun{g}() \xrightarrow{\beta} \fun{g}()
  \end{equation*}
  We have \(\fun{f}(\fun{g}()) \xrightarrow{\alpha} 0\) but
  \(\fun{f}(\fun{g}()) \xrightarrow{\beta} \fun{f}(\fun{g}())
  \xrightarrow{\beta} \dots\)

  \bigskip

  Despite the loss of expressiveness, we shall retain in this lecture
  the call by value, as it is the strategy of \OCaml and it is easy to
  follow.

  \bigskip

  In general, the order in which the arguments of a function call are
  evaluated is not specified.

\end{frame}

% ------------------------------------------------------------------------
%
\begin{frame}{Call-by-value}

  Another reason to choose call-by-value is that it allows us to
  restrict the shape of the left\hyp{}hand sides, called
  \textbf{patterns}, to one, outermost function call.

  \bigskip

  For instance, we can then disallow, for being useless, a
  rule like

  \begin{equation*}
    \fun{plus}(x,\fun{plus}(y,z)) \rightarrow
    \fun{plus}(\fun{plus}(x,y),z).
  \end{equation*}

  \noindent because, following call\hyp{}by\hyp{}value, the argument
  \(\fun{plus}(y,z)\) must be evaluated before the outermost call,
  therefore it cannot be part of any pattern.

\end{frame}

% ------------------------------------------------------------------------
%
\begin{frame}{Termination}

  We do not require by construction that all functional programs
  terminate, although that is a desirable property, but we
  nevertheless speak of `functions' instead of `partial functions'
  (which would be technically correct).

  \bigskip

  If we enforce termination, we loose
  \textbf{Turing\hyp{}completeness}: we cannot program all computable
  functions.

  \bigskip

  The termination of rewrite systems and, in particular, functional
  programs, is an \textbf{undecidable problem} in general.

  \bigskip

  Nevertheless, there exists a set of well\hyp{}known \textbf{decision
    criteria}, which we will use when programming recursive functions
  in \OCaml.

\end{frame}

% ------------------------------------------------------------------------
%
\begin{frame}{Higher-order programs, recursion, $\lambda$-calculus}

  Most functional languages allow \textbf{higher\hyp{}order function}
  definitions, whereas standard term\hyp{}rewriting systems do
  not. For example:
  \begin{equation*}
    \fun{f}(g,0) \rightarrow 1
    \qquad
    \fun{f}(g,n) \rightarrow n \times g(g,n-1)
    \qquad
    \fun{fact}(n) \rightarrow \fun{f}(\fun{f},n)
  \end{equation*}
  where \(n \in \mathbb{N}\). Note that these two definitions are
  \textbf{not} recursive, yet the function \fun{fact} is the factorial
  function. A function definition is \textbf{recursive} if the name of
  the definition occurs on at least one right\hyp{}hand side.

  \bigskip

  An adequate theoretical framework to understand higher\hyp{}order
  functions is \textbf{\(\lambda\)-calculus}. In fact,
  \(\lambda\)-calculus features prominently in the semantics of
  programming languages, even not functional ones.

  \bigskip

  For didactic purposes, we began with rewrite systems because they
  offer implicit pattern matching (that is, rule selection).

\end{frame}

% ------------------------------------------------------------------------
%
\begin{frame}{Stacks in a functional setting}

  We show how to express linear data structures in a purely functional
  language and how to compute with them. The simplest of those
  structures is the \textbf{stack}.

  \bigskip

  A stack can be thought of as a finite series of items that can only
  be accessed sequentially from one end, called the \textbf{top},
  following the analogy with a stack of material objects.

\end{frame}

% ------------------------------------------------------------------------
%
\begin{frame}{Stacks in a functional setting}

  To define a stack in a functional language, one uses calls to
  undefined functions to model all the possible shapes a stack can
  take.

  \bigskip

  Here, a stack can be either empty or not.

  \bigskip

  Let \(\fun{nil}()\) be the irreducible call that denotes an empty
  stack (that is, the function `\fun{nil}' is undefined).

  \bigskip

  If not empty, the stack then must contain a topmost (first) item,
  and a \textbf{substack}, that is, the stack made of all the items
  but the first.

  \bigskip

  Let \(\fun{cons}(x,s)\) denote the non\hyp{}empty stack with the
  item~\(x\) on top and the \textbf{immediate substack}~\(s\).

\end{frame}

% ------------------------------------------------------------------------
%
\begin{frame}{Inductive definition of stacks}

  \begin{itemize}

    \item The previous definition is intuitive, but not formal yet.

    \item Data structures are usually defined \textbf{inductively}.

    \item Let~\(\mathcal{T}\) be the set of all possible terms and
      \(\mathcal{S} \subset \mathcal{T}\) be the set of all stacks.

    \item Formally, \(\mathcal{S}\)~can by defined by
      \textbf{induction} as the smallest set such that
      \begin{enumerate}

        \item \(\fun{nil}() \in \mathcal{S}\);

        \item if \(x \in \mathcal{T}\) and \(s \in \mathcal{S}\), then
          \(\fun{cons}(x,s) \in \mathcal{S}\).
      \end{enumerate}

    \item We take the smallest set because we do not want
      in~\(\mathcal{S}\) terms that were not constructed by the
      inductive rules (this is called the \textbf{smallest fixed
        point} of the inductive relation).

  \end{itemize}

\end{frame}

% ------------------------------------------------------------------------
%
\begin{frame}{Appending a stack}

  \begin{itemize}

    \item For all \(s, t \in \mathcal{S}\), let us consider the
      functional program
      \begin{align*}
        \fun{cat}(\fun{nil}(),t) &\xrightarrow{\alpha} t\\
        \fun{cat}(\fun{cons}(x,s),t) &\xrightarrow{\smash\beta}
        \fun{cons}(x,\fun{cat}(s,t))
      \end{align*}

    \item This is a functional program suitable for
      call\hyp{}by\hyp{}value reduction strategy, because the function
      calls in the patterns (that is, left\hyp{}hand sides) are
      irreducible.

    \item The right\hyp{}hand sides are stacks, so a call to
      `\fun{cat}' is a stack.

    \item Rule~\(\alpha\) says that, if the first stack is empty, the
      value is the second.

    \item In rule~\(\beta\), the first stack is not empty, and its top
      item is the top of the final stack~(\(x\)), the immediate
      substack~(\(s\)) being used in a call to `\fun{cat}' with the
      second stack~(\(t\)) unchanged.

  \end{itemize}

\end{frame}

% ------------------------------------------------------------------------
%
\begin{frame}{Appending a stack}

  This function appends the second stack at the bottom of a
  \textbf{copy} of the first. (The function name `\fun{cat}' stands
  for `catenation', and is slightly better than `\fun{app}').

  \bigskip

  We could also say that this function appends a copy of the first
  stack on top of the other. But this is not the same as pushing the
  first stack on top of the other, like
  \begin{equation*}
    \fun{cat}(s,t) \rightarrow \fun{cons}(s,t) \qquad
    \text{Wrong!}
  \end{equation*}

  This is \textbf{not} appending: here, the stack~\(s\) becomes an
  item of the resulting stack, but we want only the \textbf{contents}
  of~\(s\) to be pushed on top of~\(t\), not their container.

  \bigskip

  Analogy: You may think of an empty stack as a empty cardboard box,
  open topside, and every item pushed in it is like a book that you
  put on the topmost book, or at the bottom of the box if this is the
  first book.

\end{frame}

% ------------------------------------------------------------------------
%
\begin{frame}{Appending a stack}

  Let us set the abbreviations
  \begin{itemize}

    \item \(\el := \fun{nil}()\),

    \item \(\cons{x}{s} := \fun{cons}(x,s)\),

  \end{itemize}
  after the convention of the programming language \OCaml.

  \bigskip

  We can further abbreviate
  \begin{itemize}

    \item \(\cons{x_1}{\cons{x_2}{\cons{\dots}{\cons{x_n}{s}}}} :=
      \fun{cons}(x_1, \fun{cons}(x_2, \dots \fun{cons}(x_n,s)))\)

    \item \([x] := \cons{x}{\el}\)

    \item \([x_1; x_2, \dots; x_n] :=
          \cons{x_1}{\cons{x_2}{\cons{x_3}{\cons{\ldots}{\cons{x_n}{\el}}}}}\).

  \end{itemize}

  Our system now becomes a bit more legible:
  \begin{align*}
    \fun{cat}(        \el,t) &\xrightarrow{\alpha} t\\
    \fun{cat}(\cons{x}{s},t) &\xrightarrow{\smash{\beta}}
    \cons{x}{\fun{cat}(s,t)}
  \end{align*}

\end{frame}

% ------------------------------------------------------------------------
%
\begin{frame}{Pattern matching}

  \begin{itemize}

    \item Let us compute \(\fun{cat}([1;2],[3])\).

    \item We need to find the first rule that is matched by this call.

    \item Remember that the arguments must be values before you do
      that.

    \item Intuitively, this means that we compare the call with the
      patterns, that is, the left\hyp{}hand sides of the rules, in
      order.

    \item Here, rule~\(\alpha\) is not matched because \([1;2] \neq
      []\).

    \item Rule~\(\beta\) is matched because \([1;2] = 1::[2]\) and
      equals the sub\hyp{}pattern \(x::s\) if we assign \(x := 1\) and
      \(s := [2]\). Also, trivially, \(t := [3]\) works.

    \item After the rule~\(\beta\) is selected, the assignments are
      used to replace the variables in the right\hyp{}hand sides.

    \item The whole evaluation is
    \begin{equation*}
      \fun{cat}([1;2],[3])
      \xrightarrow{\beta}
      \cons{1}{\fun{cat}([2],[3])}
      \xrightarrow{\beta}
      \cons{1}{\cons{2}{\fun{cat}(\el,[3])}}
      \xrightarrow{\alpha}
                    [1;2;3].
    \end{equation*}
  \end{itemize}

\end{frame}

% ------------------------------------------------------------------------
%
\begin{frame}{Abstract Syntax Trees (ASTs)}

  \begin{itemize}

    \item Depending on the context, we may use the arborescent
      depiction of terms to bring to the fore certain aspects of a
      computation.

    \item For example, it may be interesting to show how parts of the
      output (the right\hyp{}hand side) are actually \textbf{shared}
      with the input (the left\hyp{}hand side).

    \item In other words, how much (new) data is created by a given
      rule.

    \item This concept supposes that terms reside in some sort of
      space and that they can be referred to from different other
      terms.

    \item This abstract space serves as a model of an interpreter's
      memory, called the \textbf{heap}.

  \end{itemize}

\end{frame}

% ------------------------------------------------------------------------
%
\begin{frame}{Data sharing}

  \begin{itemize}

    \item Consider for instance
      \begin{center}
        \includegraphics[bb=71 656 302 721]{cat_dag}
      \end{center}
      which is the same definition of `\fun{cat}' as given above:
      \begin{equation*}
        \fun{cat}(\el,t) \xrightarrow{\alpha} t
        \qquad\qquad\qquad
        \fun{cat}(\cons{x}{s},t) \xrightarrow{\beta}
        \cons{x}{\fun{cat}(s,t)}
      \end{equation*}

    \item The arrows on certain edges denote data sharing.

    \item When trees are used to visualise terms, they are called
      \textbf{abstract syntax trees} (AST).

    \item When some trees share subtrees, the whole \textbf{forest}
      (set of trees) is called a \textbf{directed acyclic graph}
      (DAG).

  \end{itemize}

\end{frame}

% ------------------------------------------------------------------------
%
\begin{frame}{Data sharing and evaluation}

  The evaluation of \(\fun{cat}([1;2],[3])\) is:
  \bigskip
  \begin{center}
    \includegraphics[bb=71 621 276 733]{cat123_push}
  \end{center}

  Note that each right-hand side shows only the rewritten call, not
  the whole term.

\end{frame}

% ------------------------------------------------------------------------
%
\begin{frame}{Data sharing and evaluation}

  Therefore, to reconstruct the value, we need to work leftwards and
  replace the rewritten call by the right-hand side:
  \begin{center}
    \includegraphics[bb=71 623 248 717]{cat123_pop0}
  \end{center}

  We do not show the calls that have been replaced by their values.

\end{frame}

% ------------------------------------------------------------------------
%
\begin{frame}{Data sharing and evaluation}

  \begin{center}
    \includegraphics[bb=71 621 212 715]{cat123_pop1}
    \qquad\quad
    \includegraphics[bb=71 621 212 715]{cat123_pop2}
  \end{center}

\end{frame}

% ------------------------------------------------------------------------
%
\begin{frame}{Call stack}

  We can see that the value of the second argument, [3], is shared
  with the value of the call, as well as the integers~1 and~2.

  \bigskip

  We saw how rewrites accumulated rightwards until the result had to
  be reconstructed by replacing leftwards the right-hand sides.

  \bigskip

  This dynamic is that of a stack, but a stack holding suspended
  function calls (that is, awaiting their values) and references to
  ASTs in the heap. This is the \textbf{call stack}.

  \bigskip

  The top of the stack (here, the rightmost abstract syntax tree) is
  the next term to be rewritten, unless the call on the left\hyp{}hand
  side is the original call, in which case the top is its value.

\end{frame}

% ------------------------------------------------------------------------
%
\begin{frame}{Garbage collection}

  \begin{itemize}

    \item In the abstract syntax trees (AST), the nodes with the calls
      that are replaced by their values have been immediately
      removed. In fact, they could have been removed later, in any
      order.

    \item The process which identifies and discards useless ASTs is
      called the \textbf{garbage collector} (GC).

    \item A term is useless if it cannot be accessed from anywhere in
      the call stack.

    \item Contrary to the heap, which can be thought of as an
      unordered set of trees (a forest of DAGs), the call stack is an
      ordered set that can be efficiently extended or reduced.

    \item The garbage collector can be conceived as a process
      interleaved with the process of evaluation.

    \item Like evaluators of functional programs, garbage collectors
      may also have different strategies (as to when dispose of
      useless data).
  \end{itemize}

\end{frame}

% ------------------------------------------------------------------------
%
\begin{frame}{Tail call optimisation}

  Let us define a function that sums integers in a non\hyp{}empty
  list:
  \begin{equation*}
    \fun{sum}([n]) \xrightarrow{\alpha} n\qquad
    \fun{sum}(\cons{n}{s}) \xrightarrow{\beta} n + \fun{sum}(s)
  \end{equation*}

  The view of the system with maximum sharing is:
  \begin{center}
    \includegraphics{sum_dag}
  \end{center}

\end{frame}

% ------------------------------------------------------------------------
%
\begin{frame}{Tail call optimisation}

  The first phase of the evaluation of \(\fun{sum}([3;7;5])\) is:
  \begin{center}
    \includegraphics[bb=72 635 404 721,scale=0.99]{sum375_push}
  \end{center}

\end{frame}

% ------------------------------------------------------------------------
%
\begin{frame}{Tail call optimisation}

  The second phase (replacing the pending calls with their values) is
  \begin{center}
    \includegraphics[bb=72 644 380 721]{sum375_pop}
  \end{center}

\end{frame}

% ------------------------------------------------------------------------
%
\begin{frame}{Tail call optimisation}

  The same, with a garbage collector that collects after each
  replacement:
  \begin{center}
    \includegraphics[bb=149 585 448 662]{sum375_pop1}
  \end{center}

\end{frame}

% ------------------------------------------------------------------------
%
\begin{frame}{Tail call optimisation}

  Let us consider an alternate version of \fun{sum}:
  \begin{equation*}
    \begin{array}{@{}r@{\;}l@{\;}l@{}}
      \fun{sum}_1(\cons{n}{t}) & \xrightarrow{\gamma}
      & \fun{sum}_0(n,t).\\
      \fun{sum}_0(n,\el) & \xrightarrow{\delta} & n\\
      \fun{sum}_0(n,\cons{m}{s}) & \xrightarrow{\epsilon}
      & \fun{sum}_0(n+m,s).
    \end{array}
  \end{equation*}

  For \(\fun{sum}_1([3;7;5])\) we have the evaluation
  \begin{center}
    \includegraphics[bb=71 641 342 722]{sum0_375_0}
  \end{center}

\end{frame}

% ------------------------------------------------------------------------
%
\begin{frame}{Tail call optimisation}

  \begin{center}
    \includegraphics[bb=71 641 346 721]{sum0_375_1}
    \includegraphics[bb=71 639 384 721]{sum0_375_2}
  \end{center}

\end{frame}

% ------------------------------------------------------------------------
%
\begin{frame}{Tail call optimisation}

  Notice that, in the case of \fun{sum}$_0$, we found the result (15)
  at the end of the first phase (of growing the call stack).

  \bigskip

  Therefore, there is no need for a second phase, which replaces the
  pending calls by their values: we can discard all those calls at
  once after the first phase.

  \bigskip

  Or, to see it even more clearly, let us rewrite the same first phase
  and assume that the collector reclaims immediately the useless
  pending calls:
  \begin{center}
    \includegraphics[bb=71 640 326 721]{sum0_375tco0}
  \end{center}

\end{frame}

% ------------------------------------------------------------------------
%
\begin{frame}{Tail call optimisation}

  \begin{center}
    \includegraphics[bb=71 628 350 721]{sum0_375tco1}
    \includegraphics[bb=71 639 366 721]{sum0_375tco2}
  \end{center}

\end{frame}

% ------------------------------------------------------------------------
%
\begin{frame}{Tail call optimisation}

  As we see now, we can keep the size of the call stack constant, if
  we reclaim the top at each new rewrite.

  \bigskip

  This optimisation is called \textbf{tail call optimisation}.

  \bigskip

  Intuitively, a call in \textbf{tail position} is a call whose value
  is the value of the previously rewritten call.

  \bigskip

  For example, the call \(f(x)\) in \(1 + f(x)\) is not in tail
  position, because its value needs to be incremented.

  \bigskip

  Most compilers of functional languages detect and implement that
  optimisation that saves memory, including the \Clang compiler of
  \textsf{GCC}.

\end{frame}



\end{document}
